% Elemento y átomo
Aunque etimológicamente la palabra átomo significa "sin corte," en el presente siglo se ha establecido que los átomos están constituidos por electrones (partículas con carga eléctrica negativa que giran en órbitas específicas alrededor del núcleo), y un núcleo el cual está a su vez formado por protones (partículas con carga eléctrica positiva) y neutrones (partículas con carga eléctrica neutra).

Los átomos son generalmente neutros porque el número de electrones que giran alrededor del núcleo con sus cargas negativas es igual al número de protones en el núcleo con sus cargas positivas.

% Isótopos
El número de protones determina la identidad del elemento químico, cuando un "mismo" elemento tiene diferentes números de neutrones, en estos casos se habla de isótopos del elemento.

A cada elemento se le asocia con el "número de orden" que le corresponde dentro de la tabla periódica, a este número se le conoce como número atómico que corresponde al número de protones que posee el núcleo, y se representa con la letra $Z$.

A la suma de protones y neutrones se le llama número de masa y se representa con la letra $A$.

% Notación de isótopos
Para diferenciar un isótopo de otros se usa la siguiente notación: 
\[
^{A}_{Z}X_{N}
\]
donde: \\
$X$ = Símbolo químico \\
$A$ = Número másico - suma de protones y neutrones \\
$Z$ = Número atómico \\
$N$ = Número de neutrones

Por ejemplo: 
\[
^{2}_{1}\text{H}_{1}
\]

% Radiactividad y radiación
\chapter{RADIATIVIDAD Y RADIACIÓN}

A medida que los núcleos se hacen más complejos por el número de partículas que contienen, el tamaño de estos incrementa y por ende las fuerzas nucleares se van debilitando mientras que las electrostáticas se incrementan, lo que ocasiona que los núcleos pierdan estabilidad, y es por esta razón (aunque no es la única) que la mayor parte de los núcleos pesados presentan radiactividad.

La radiactividad se manifiesta a través de la emisión espontánea de ciertas radiaciones por parte de los núcleos atómicos. Dichas radiaciones son capaces de impresionar placas fotográficas.
